\documentclass[10pt,a4paper]{amsart}
\usepackage[margin=19mm]{geometry}
\usepackage{amsmath,amssymb,mathtools}
\usepackage[scr=rsfs,cal=euler]{mathalfa}
\usepackage{xfrac}
\usepackage[unicode,hidelinks,bookmarks=false]{hyperref}
\newcommand{\ie}{i.e.\ }
\newcommand{\cf}{cf.\ }
\newcommand{\resp}{respectively\ }
\newcommand{\Subsection}{\subsection}
\providecommand{\nobreakdash}{\nobreak\hbox{-}}
\setlength{\emergencystretch}{2em}
\multlinegap0pt
\usepackage[shortlabels]{enumitem}
\usepackage{xcolor}
\usepackage{amssymb}
\usepackage{algorithm}\usepackage{algpseudocode}
\newtheorem{assumption}{Assumption}
\newtheorem{proposition}{Proposition}
\newtheorem{lemma}{Lemma}
\newcommand{\proofstep}[1]{\par\smallskip\noindent\textbf{#1}\ }
\title{A Barrier-Metric First-Order Method for Linearly Constrained Bilevel Optimization}
\author{Tenglong Hong, Paul Grigas}
\date{Source: arXiv:2605.11476v1 (CC BY 4.0)}
\begin{document}
\maketitle

\noindent\emph{Source and licence.} A Barrier-Metric First-Order Method for Linearly Constrained Bilevel Optimization. Version arXiv:2605.11476v1 (CC BY 4.0), distributed by arXiv under the Creative Commons Attribution 4.0 International licence, http://creativecommons.org/licenses/by/4.0/, as recorded in the arXiv licence metadata for this exact version and shown on the abstract page https://arxiv.org/abs/2605.11476v1. Full licence notice: \texttt{licenses/arxiv-2605.11476-CC-BY-4.0.txt}.\par\smallskip

\noindent\textit{Editorial scope.} Counted unit: the article's lemma lem:proxy\_direction\_linearization (source0.tex lines 6671--6696). The article's complete original proof of it is delivered with it (source0.tex lines 6697--7010, one \texttt{proof} environment of 314 lines). No mathematics is written, repaired or re-derived: the statement and the proof are the article's own lines, in the article's own notation, and no retained line is substituted or re-flowed.  Retained uncounted as the source-local context the proof needs: lines 202--216; lines 423--460.  Nothing was omitted from any retained span, so the package carries no omission record.  No retained line contains a citation, so the wrapper carries no bibliography and no bibliography entry.  No retained line contains a figure environment, an image-inclusion command or drawing code, so no image is delivered and the package declares no asset dependency.  Editorial formatting only, in addition: an AMS \texttt{amsart} preamble that restates the article's own declarations and macros used by the retained lines, namely the environments \texttt{assumption}, \texttt{proposition}, \texttt{lemma} on separate counters, together with the standard packages those lines need.  The retained environments print in this document as Assumption 1, Proposition 1, Lemma 1 in order of appearance, whereas the article prints the same items with the numbering of its own full document.  The complete native source of the article as distributed in the arXiv e-print for this exact version is bundled unchanged as \texttt{sources/200361/source0.tex}.
\begin{assumption}[Objective regularity]
\label{ass:euclidean}
\leavevmode
\begin{enumerate}[leftmargin=*, label=\textup{(A\arabic*)}, ref=\textup{(A\arabic*)}]
    \item\label{ass:A1} \(f\) and \(g\) are jointly smooth in \((x,y)\), with constants
    \(\ell_{f,1}\) and \(\ell_{g,1}\).
    \item\label{ass:A2} For every \(x\), the map \(y\mapsto g(x,y)\) is \(\rho_g\)-strongly convex for some \(\rho_g>0\).
    \item\label{ass:A3} \(f,g\in C^2\), and \(\nabla^2 f,\nabla^2 g\) are jointly Lipschitz in
    \((x,y)\), with constants \(\ell_{f,2}\) and \(\ell_{g,2}\).
    \item\label{ass:A4} The first derivatives used in lower-level and outer analysis are bounded: for all
    \((x,y)\),
    \(\|\nabla_x f(x,y)\|_2\vee \|\nabla_y f(x,y)\|_2\le \ell_{f,0}\),
    \(\|\nabla_x g(x,y)\|_2\le \ell_{g,0}\).
\end{enumerate}
\end{assumption}

\begin{proposition}[Local Dikin regularity]
\label{prop:derived_dikin_regularity}
Suppose Assumption~\ref{ass:euclidean} holds. Fix a neighborhood radius
\(\eta\in(0,1/2)\). Then there exist finite constants
\(\rho_{\psi}^{\eta},\rho_{\psi,\lambda}^{\eta}>0\) and
\(\ell_{\psi,1}^{\eta},\ell_{\psi,2}^{\eta},
\ell_{\psi,\lambda,1}^{\eta},\ell_{\psi,\lambda,2}^{\eta},
\ell_{f,0}^{\eta},\ell_{f,1}^{\eta},\ell_{f,2}^{\eta}<\infty\),
such that, for every \(x\in\mathbb R^{d_x}\) the following hold.
\begin{enumerate}[leftmargin=*]
\item[\textup{(i)}]
On the neighborhood \(T_{\eta,\mu}(x)\), the barrierized lower objective is locally well conditioned in
the anchored Dikin metric:
\(\rho_{\psi}^{\eta}H_\mu^\star(x)\preceq
\nabla_{yy}^2\psi_\mu(x,y)\preceq
\ell_{\psi,1}^{\eta}H_\mu^\star(x)\) for every \(y\in T_{\eta,\mu}(x)\).
Moreover, \(\nabla_{xy}^2\psi_\mu\) is locally bounded, and
\(\nabla_{yy}^2\psi_\mu,\nabla_{xy}^2\psi_\mu\) are locally Lipschitz on
\(T_{\eta,\mu}(x)\) in the anchored Dikin norms, with constants
\(\ell_{\psi,1}^{\eta},\ell_{\psi,2}^{\eta}\).

\item[\textup{(ii)}]
The upper-level derivatives of \(f\) are locally controlled on \(T_{\eta,\mu}(x)\):
\(\nabla_y f\), \(\nabla_{xy}^2 f\), and \(\nabla_{yy}^2 f\) are locally bounded, and the
corresponding Hessian terms are locally Lipschitz, in the anchored Dikin norms, with constants
\(\ell_{f,0}^{\eta},\ell_{f,1}^{\eta},\ell_{f,2}^{\eta}\).

\item[\textup{(iii)}]
On the proxy neighborhood \(T_{\eta,\mu}^{\lambda}(x)\), the barrierized lower objective inherits the same
type of local Dikin regularity with respect to the proxy anchor \(H_{\lambda,\mu}^\star(x)\):
\(\rho_{\psi,\lambda}^{\eta}H_{\lambda,\mu}^\star(x)\preceq
\nabla_{yy}^2\psi_\mu(x,y)\preceq
\ell_{\psi,\lambda,1}^{\eta}H_{\lambda,\mu}^\star(x)\) for every
\(y\in T_{\eta,\mu}^{\lambda}(x)\), together with the analogous mixed-Hessian and
Lipschitz-Hessian bounds in the proxy anchored Dikin norm, with constants
\(\ell_{\psi,\lambda,1}^{\eta},\ell_{\psi,\lambda,2}^{\eta}\).
\end{enumerate}
\end{proposition}

\begin{lemma}[Local proxy-gradient linearization]
\label{lem:proxy_direction_linearization}
Fix \(\lambda>0\) and \(x\in X\). For any \(y\in T_{\eta,\mu}(x)\), we have
\begin{align}
&\Big\|
\nabla F_\mu(x)
-\nabla_x L_{\lambda,\mu}(x,y)
+\nabla^2_{xy}\psi_\mu(x,y_\mu^\star(x))^\top
\big(\nabla^2_{yy}\psi_\mu(x,y_\mu^\star(x))\big)^{-1}
\nabla_y L_{\lambda,\mu}(x,y)
\Big\|_2
\nonumber\\
&\qquad\le
2\Big(\frac{\ell_{\psi,1}^\eta}{\rho_\psi^\eta}\Big)
\|y-y_\mu^\star(x)\|_{y_\mu^\star(x)}
\Big(
\ell_{f,1}^\eta
+
\lambda\min\{
2\ell_{\psi,1}^\eta,\,
\ell_{\psi,2}^\eta\|y-y_\mu^\star(x)\|_{y_\mu^\star(x)}
\}
\Big).
\label{eq:proxy_linearization_bound}
\end{align}
\end{lemma}

\begin{proof}
Fix \(x\in X\). Since
\(\nabla_y\psi_\mu(x,y_\mu^\star(x))=0\) and
\(\nabla^2_{yy}\psi_\mu(x,y_\mu^\star(x))\succ0\) on the maintained tube, the implicit function
theorem gives
\begin{equation}
\label{eq:support_Dy_mu}
Dy_\mu^\star(x)
=
-\big(\nabla^2_{yy}\psi_\mu(x,y_\mu^\star(x))\big)^{-1}
\nabla^2_{xy}\psi_\mu(x,y_\mu^\star(x)).
\end{equation}
Hence, by the chain rule,
\begin{align}
\nabla F_\mu(x)
&=
\nabla_x f(x,y_\mu^\star(x))
+
\big(Dy_\mu^\star(x)\big)^\top
\nabla_y f(x,y_\mu^\star(x))
\nonumber\\
&=
\nabla_x f(x,y_\mu^\star(x))
-
\nabla^2_{xy}\psi_\mu(x,y_\mu^\star(x))^\top
\big(\nabla^2_{yy}\psi_\mu(x,y_\mu^\star(x))\big)^{-1}
\nabla_y f(x,y_\mu^\star(x)).
\label{eq:support_gradF_mu}
\end{align}

Next, using
\(L_{\lambda,\mu}(x,y)=f(x,y)+\lambda(\psi_\mu(x,y)-\psi_\mu^\star(x))\) and
\(\nabla_x\psi_\mu^\star(x)=\nabla_x\psi_\mu(x,y_\mu^\star(x))\), we have
\begin{align}
\nabla_x L_{\lambda,\mu}(x,y)
&=
\nabla_x f(x,y)
+
\lambda\big(
\nabla_x\psi_\mu(x,y)-\nabla_x\psi_\mu(x,y_\mu^\star(x))
\big),
\nonumber\\
\nabla_y L_{\lambda,\mu}(x,y)
&=
\nabla_y f(x,y)+\lambda\nabla_y\psi_\mu(x,y).
\label{eq:support_proxy_derivatives}
\end{align}

Combining \eqref{eq:support_gradF_mu}--\eqref{eq:support_proxy_derivatives} gives
\begin{align}
&\nabla F_\mu(x)
-\nabla_x L_{\lambda,\mu}(x,y)
+
\nabla^2_{xy}\psi_\mu(x,y_\mu^\star(x))^\top
\big(\nabla^2_{yy}\psi_\mu(x,y_\mu^\star(x))\big)^{-1}
\nabla_y L_{\lambda,\mu}(x,y)
\nonumber\\
&=
\underbrace{
\nabla_x f(x,y_\mu^\star(x))-\nabla_x f(x,y)
}_{T_1}
\nonumber\\
&\quad+
\underbrace{
\nabla^2_{xy}\psi_\mu(x,y_\mu^\star(x))^\top
\big(\nabla^2_{yy}\psi_\mu(x,y_\mu^\star(x))\big)^{-1}
\big(
\nabla_y f(x,y)-\nabla_y f(x,y_\mu^\star(x))
\big)
}_{T_2}
\nonumber\\
&\quad
-\lambda
\underbrace{
\Big(
\nabla_x\psi_\mu(x,y)-\nabla_x\psi_\mu(x,y_\mu^\star(x))
-\nabla^2_{xy}\psi_\mu(x,y_\mu^\star(x))^\top
(y-y_\mu^\star(x))
\Big)
}_{R_x}
\nonumber\\
&\quad
+\lambda
\underbrace{
\nabla^2_{xy}\psi_\mu(x,y_\mu^\star(x))^\top
\big(\nabla^2_{yy}\psi_\mu(x,y_\mu^\star(x))\big)^{-1}
}_{\text{linear correction}}
\nonumber\\
&\qquad\qquad\qquad\cdot
\underbrace{
\Big(
\nabla_y\psi_\mu(x,y)-\nabla_y\psi_\mu(x,y_\mu^\star(x))
-\nabla^2_{yy}\psi_\mu(x,y_\mu^\star(x))(y-y_\mu^\star(x))
\Big)
}_{\widetilde R_y}.
\label{eq:support_proxy_decomposition}
\end{align}
The linear \(\lambda\nabla^2_{xy}\psi_\mu(x,y_\mu^\star(x))^\top(y-y_\mu^\star(x))\) terms cancel.

\proofstep{Operator-norm in the anchored Dikin geometry:}
We will use the induced norms
\[
\|M\|_{2\to y_\mu^\star(x),*}
:=
\sup_{\|u\|_2=1}\|Mu\|_{y_\mu^\star(x),*},
\qquad
\|M\|_{y_\mu^\star(x)\to 2}
:=
\sup_{\|v\|_{y_\mu^\star(x)}=1}\|Mv\|_2.
\]
By duality,
\[
\|M^\top\|_{y_\mu^\star(x)\to 2}
=
\|M\|_{2\to y_\mu^\star(x),*}.
\]
Proposition~\ref{prop:derived_dikin_regularity} implies
\begin{equation}
\label{eq:support_cross_operator_bound}
\|\nabla^2_{xy}\psi_\mu(x,y_\mu^\star(x))\|_{2\to y_\mu^\star(x),*}
\le
\ell_{\psi,1}^\eta,
\end{equation}
and therefore
\[
\|\nabla^2_{xy}\psi_\mu(x,y_\mu^\star(x))^\top\|_{y_\mu^\star(x)\to 2}
\le
\ell_{\psi,1}^\eta .
\]
Also, since
\[
\nabla^2_{yy}\psi_\mu(x,y_\mu^\star(x))
\succeq
\rho_\psi^\eta H_\mu^\star(x),
\]
we have
\begin{equation}
\label{eq:support_hinv_operator_bound}
\big\|
\big(\nabla^2_{yy}\psi_\mu(x,y_\mu^\star(x))\big)^{-1}w
\big\|_{y_\mu^\star(x)}
\le
\frac{1}{\rho_\psi^\eta}\|w\|_{y_\mu^\star(x),*}.
\end{equation}
Combining the preceding two displays yields
\begin{equation}
\label{eq:support_mixed_operator_bound}
\Big\|
\nabla^2_{xy}\psi_\mu(x,y_\mu^\star(x))^\top
\big(\nabla^2_{yy}\psi_\mu(x,y_\mu^\star(x))\big)^{-1}
\Big\|_{y_\mu^\star(x),*\to 2}
\le
\frac{\ell_{\psi,1}^\eta}{\rho_\psi^\eta}.
\end{equation}

\proofstep{Bound \(T_1\) and \(T_2\).}
By the local bound for \(f\), with Euclidean--Dikin conversion factors absorbed into
\(\ell_{f,1}^\eta\),
\begin{equation}
\label{eq:support_T1_bound}
\|T_1\|_2
\le
\ell_{f,1}^\eta\|y-y_\mu^\star(x)\|_{y_\mu^\star(x)}.
\end{equation}
For \(T_2\), combine \eqref{eq:support_mixed_operator_bound} with the local Lipschitz bound for
\(\nabla_y f\):
\[
\|T_2\|_2
\le
\frac{\ell_{\psi,1}^\eta}{\rho_\psi^\eta}
\|\nabla_y f(x,y)-\nabla_y f(x,y_\mu^\star(x))\|_{y_\mu^\star(x),*}
\le
\frac{\ell_{\psi,1}^\eta}{\rho_\psi^\eta}
\ell_{f,1}^\eta
\|y-y_\mu^\star(x)\|_{y_\mu^\star(x)} .
\]
Since \(\ell_{\psi,1}^\eta/\rho_\psi^\eta\ge1\), \eqref{eq:support_T1_bound} gives
\begin{equation}
\label{eq:support_f_terms_bound}
\|T_1\|_2+\|T_2\|_2
\le
2\Big(\frac{\ell_{\psi,1}^\eta}{\rho_\psi^\eta}\Big)
\ell_{f,1}^\eta
\|y-y_\mu^\star(x)\|_{y_\mu^\star(x)} .
\end{equation}

\proofstep{Bound \(R_x\) and \(\widetilde R_y\).}
For \(R_x\), the integral remainder form gives
\begin{align}
R_x
&=
\int_0^1
\Big(
\nabla^2_{xy}\psi_\mu(x,y_\mu^\star(x)+t(y-y_\mu^\star(x)))
-
\nabla^2_{xy}\psi_\mu(x,y_\mu^\star(x))
\Big)^\top
\nonumber\\
&\hspace{6.5cm}\cdot
(y-y_\mu^\star(x))\,dt .
\label{eq:support_Rx_integral}
\end{align}
By duality, for any matrix \(M\),
\[
\|M^\top(y-y_\mu^\star(x))\|_2
\le
\|M\|_{2\to y_\mu^\star(x),*}
\|y-y_\mu^\star(x)\|_{y_\mu^\star(x)}.
\]
Thus
\begin{align}
\|R_x\|_2
&\le
\|y-y_\mu^\star(x)\|_{y_\mu^\star(x)}
\int_0^1
\Big\|
\nabla^2_{xy}\psi_\mu(x,y_\mu^\star(x)+t(y-y_\mu^\star(x)))
\nonumber\\
&\hspace{4.5cm}
-
\nabla^2_{xy}\psi_\mu(x,y_\mu^\star(x))
\Big\|_{2\to y_\mu^\star(x),*}\,dt .
\label{eq:support_Rx_prebound}
\end{align}
On the tube, either use the crude bound
\(\|\nabla^2_{xy}\psi_\mu(x,\cdot)\|_{2\to y_\mu^\star(x),*}\le\ell_{\psi,1}^\eta\) twice, or use the
Hessian-Lipschitz-in-\(y\) bound with constant \(\ell_{\psi,2}^\eta\). Hence
\begin{equation}
\label{eq:support_Rx_bound}
\|R_x\|_2
\le
\|y-y_\mu^\star(x)\|_{y_\mu^\star(x)}
\min\Big\{
2\ell_{\psi,1}^\eta,\,
\ell_{\psi,2}^\eta\|y-y_\mu^\star(x)\|_{y_\mu^\star(x)}
\Big\}.
\end{equation}

Similarly, since \(\nabla_y\psi_\mu(x,y_\mu^\star(x))=0\),
\begin{align}
\widetilde R_y
&=
\nabla_y\psi_\mu(x,y)
-
\nabla_y\psi_\mu(x,y_\mu^\star(x))
-
\nabla^2_{yy}\psi_\mu(x,y_\mu^\star(x))(y-y_\mu^\star(x))
\nonumber\\
&=
\int_0^1
\Big(
\nabla^2_{yy}\psi_\mu(x,y_\mu^\star(x)+t(y-y_\mu^\star(x)))
-
\nabla^2_{yy}\psi_\mu(x,y_\mu^\star(x))
\Big)
\nonumber\\
&\hspace{6.5cm}\cdot
(y-y_\mu^\star(x))\,dt .
\label{eq:support_Ry_integral}
\end{align}
Using the corresponding operator norm
\(\|\cdot\|_{y_\mu^\star(x)\to y_\mu^\star(x),*}\), we get
\begin{equation}
\label{eq:support_Ry_inner_bound}
\|\widetilde R_y\|_{y_\mu^\star(x),*}
\le
\|y-y_\mu^\star(x)\|_{y_\mu^\star(x)}
\min\Big\{
2\ell_{\psi,1}^\eta,\,
\ell_{\psi,2}^\eta\|y-y_\mu^\star(x)\|_{y_\mu^\star(x)}
\Big\}.
\end{equation}
Finally, applying \eqref{eq:support_mixed_operator_bound} to the linear correction multiplying
\(\widetilde R_y\), we obtain
\begin{equation}
\label{eq:support_Ry_bound}
\Big\|
\nabla^2_{xy}\psi_\mu(x,y_\mu^\star(x))^\top
\big(\nabla^2_{yy}\psi_\mu(x,y_\mu^\star(x))\big)^{-1}
\widetilde R_y
\Big\|_2
\le
\frac{\ell_{\psi,1}^\eta}{\rho_\psi^\eta}
\|\widetilde R_y\|_{y_\mu^\star(x),*}.
\end{equation}
Since \(\ell_{\psi,1}^\eta/\rho_\psi^\eta\ge1\), combining
\eqref{eq:support_Rx_bound}--\eqref{eq:support_Ry_bound} gives
\begin{align}
&\lambda
\left(
\|R_x\|_2
+
\Big\|
\nabla^2_{xy}\psi_\mu(x,y_\mu^\star(x))^\top
\big(\nabla^2_{yy}\psi_\mu(x,y_\mu^\star(x))\big)^{-1}
\widetilde R_y
\Big\|_2
\right)
\nonumber\\
&\qquad\le
2\Big(\frac{\ell_{\psi,1}^\eta}{\rho_\psi^\eta}\Big)
\lambda
\|y-y_\mu^\star(x)\|_{y_\mu^\star(x)}
\min\Big\{
2\ell_{\psi,1}^\eta,\,
\ell_{\psi,2}^\eta\|y-y_\mu^\star(x)\|_{y_\mu^\star(x)}
\Big\}.
\label{eq:support_remainder_bound}
\end{align}

Taking Euclidean norms in \eqref{eq:support_proxy_decomposition}, applying the triangle inequality,
and using \eqref{eq:support_f_terms_bound} and \eqref{eq:support_remainder_bound}, we obtain
\eqref{eq:proxy_linearization_bound}.
\end{proof}

\end{document}
